\section*{Problems}

\subsection*{Problem Statements}

% Remember
\begin{problema}
\label{prob:ca86634}
To validate residential electricity-demand forecasts at an urban substation ($Y$, megawatt-hours), an engineer evaluates a regression model on $n_{\text{test}}=4$ independent validation days:
\begin{center}
\begin{tabular}{ccc}
\hline
\textbf{Day ($i$)} & \textbf{Actual consumption ($Y_i$)} & \textbf{Model prediction ($\hat Y_i$)} \\ \hline
1 & 120.0 & 115.0 \\
2 & 135.0 & 138.0 \\
3 & 110.0 & 112.0 \\
4 & 145.0 & 141.0 \\ \hline
\end{tabular}
\end{center}
\begin{enumerate}
\item Calculate test Mean Squared Error ($\text{MSE}_{\text{test}}=\frac1{n_{\text{test}}}\sum(Y_i-\hat Y_i)^2$).
\item Calculate test Root Mean Squared Error ($\text{RMSE}_{\text{test}}$) and Mean Absolute Error ($\text{MAE}_{\text{test}}$).
\item Explain why RMSE is always at least as large as MAE and what geometric penalty it imposes for large forecast errors.
\end{enumerate}
\end{problema}

% Understand
\begin{problema}
\label{prob:7d813ee}
A data-science team fits a multiple-regression model for residential property value ($Y$, thousands of dollars). With $n=200$ observations (80\% training, 20\% test), $R^2_{\text{train}}=0.760$, $RSE_{\text{train}}=24.50$, $R^2_{\text{test}}=0.720$, and $RSE_{\text{test}}=27.80$.
\begin{enumerate}
\item Assess evidence of overfitting by comparing in-sample and out-of-sample metric gaps.
\item If the mean test-set value is $\bar Y_{\text{test}}=250.0$ thousand dollars, calculate the mean residual predictive-error percentage relative to the market mean.
\item Explain why a simple two-way split (Holdout Method) can be sensitive or unstable for moderate samples, and propose a superior technical alternative.
\end{enumerate}
\end{problema}

% Apply
\begin{problema}
\label{prob:d0d57bc}
When auditing a linear econometric model with $p=15$ variables over $n=120$ corporate transactions, split into 80\% calibration ($n_{\text{train}}=96$) and 20\% evaluation ($n_{\text{test}}=24$), the metrics are:
\begin{center}
\begin{tabular}{|l|c|c|}\hline
\textbf{Metric} & \textbf{Calibration} & \textbf{Evaluation} \\ \hline
$R^2$ & 0.850 & 0.620 \\
$RSE$ & 15.20 & 28.70 \\
$MAE$ & 12.10 & 22.50 \\ \hline
\end{tabular}
\end{center}
\begin{enumerate}
\item Calculate out-of-sample degradation ratios $RSE_{\text{test}}/RSE_{\text{train}}$ and $R^2_{\text{test}}/R^2_{\text{train}}$.
\item Diagnose whether the model has \textbf{extreme overfitting (high variance)} or \textbf{underfitting (high bias)}, justifying the conclusion with standard thresholds.
\item Prepare a correction plan with at least four precise technical interventions (regularization, variable selection, nested cross-validation, and so on).
\end{enumerate}
\end{problema}

% Analyze
\begin{problema}
\label{prob:cb39746}
To model price elasticity of public-transport demand ($n=100$ trips), three models with different complexity are fitted ($p$ includes the intercept and $\sigma^2$):
\begin{center}
\begin{tabular}{|l|c|c|c|}\hline
\textbf{Model} & \textbf{$p$} & \textbf{SSE} & \textbf{SSE$/n$} \\ \hline
A (Basic) & 3 & 450.00 & 4.500 \\
B (Intermediate) & 5 & 380.00 & 3.800 \\
C (Saturated) & 8 & 350.00 & 3.500 \\ \hline
\end{tabular}
\end{center}
Use $\text{AIC}=n\ln(\text{SSE}/n)+2p$, $\text{BIC}=n\ln(\text{SSE}/n)+p\ln(n)$, with $\ln(4.50)=1.5041$, $\ln(3.80)=1.3350$, $\ln(3.50)=1.2528$, and $\ln(100)=4.6052$.
\begin{enumerate}
\item Manually calculate AIC and BIC for all three models.
\item Determine which model is selected under AIC and which under BIC.
\item Explain why BIC imposes an asymptotically more severe complexity penalty than AIC when $n\ge8$, and discuss implications for model selection in Big Data.
\end{enumerate}
\end{problema}

% Evaluate
\begin{problema}
\label{prob:2b2e597}
Two retail-sales predictive architectures are compared under 10-fold cross-validation:
\begin{itemize}
\item \textbf{Model A ($p_A=3$):} $\bar R^2_{\text{CV},A}=0.780\pm0.040$ (mean $\pm$ fold standard deviation).
\item \textbf{Model B ($p_B=5$):} $\bar R^2_{\text{CV},B}=0.800\pm0.060$.
\end{itemize}
\begin{enumerate}
\item Determine which model has greater mean explanatory power and which is more stable using the percentage coefficient of variation.
\item If both are deployed over $n=150$ stores, calculate the expected adjusted $R^2$ for each.
\item Applying parsimony (Occam's razor), evaluate which model management should adopt, weighing the marginal fit gain against complexity cost and greater instability of the saturated model.
\end{enumerate}
\end{problema}

% Create
\begin{problema}
\label{prob:9689172}
Design an original model-validation problem (different from property, electricity, corporate-transaction, public-transport, or retail-sales examples): a streaming platform evaluates a subscription-cancellation (\emph{churn}) model. It splits $n=500$ customers into 350 training and 150 test cases, obtaining $R^2_{\text{train}}=0.82$, $RSE_{\text{train}}=8.5$, $R^2_{\text{test}}=0.79$, and $RSE_{\text{test}}=9.6$. Calculate the $R^2$ gap and relative RSE increase, and assess whether overfitting is mild/admissible or severe/pathological under the standard scale ($\Delta R^2\le0.05$ and RSE increase $\le15\%$ admissible; $\Delta R^2>0.15$ or increase $>40\%$ severe).
\end{problema}

\subsection*{Detailed Solutions}

% Remember
\begin{solproblema}[prob:ca86634]
Residuals are $e_1=+5.0$, $e_2=-3.0$, $e_3=-2.0$, and $e_4=+4.0$, with $\sum e_i^2=54.00$. Thus $\text{MSE}_{\text{test}}=54.00/4=13.50$ (MWh)$^2$. $\text{RMSE}_{\text{test}}=\sqrt{13.50}\approx3.6742$ MWh. Since $\sum|e_i|=14.00$, $\text{MAE}_{\text{test}}=14.00/4=3.5000$ MWh. RMSE $\ge$ MAE by the inequality between the quadratic-mean ($L_2$) and absolute-mean ($L_1$) norms, with equality only when all $|e_i|$ are identical. Squaring each residual before averaging disproportionately penalizes large errors: an error of 10 MWh contributes 100 units, equivalent to one hundred errors of 1 MWh.
\end{solproblema}

% Understand
\begin{solproblema}[prob:7d813ee]
\begin{enumerate}
\item $\Delta R^2=0.760-0.720=0.040$ (4.0 points). Relative RSE increase is $(27.80-24.50)/24.50\times100\approx13.47\%$. Under the standard scale ($\Delta R^2\le0.05$ and RSE increase $\le15\%$ is mild/admissible), there is \textbf{no pathological overfitting}; the model generalizes reasonably.
\item Relative forecast error is $RSE_{\text{test}}/\bar Y_{\text{test}}\times100=27.80/250.00\times100=11.12\%$.
\item With moderate $n=200$, reserving one fixed block of 40 test observations creates high variance (the estimate depends on which observations happen to be held out) and loses calibration power. A superior alternative is $k$-fold cross-validation ($k=5$ or $k=10$), rotating partitions so each datum is used exactly once for testing and averaging the generalization error.
\end{enumerate}
\end{solproblema}

% Apply
\begin{solproblema}[prob:d0d57bc]
\begin{enumerate}
\item $RSE_{\text{test}}/RSE_{\text{train}}=28.70/15.20\approx1.8882$ (the error nearly doubles, $+88.82\%$). $R^2_{\text{test}}/R^2_{\text{train}}=0.620/0.850\approx0.7294$ (a 27.06\% loss in explained variance).
\item The model exhibits \textbf{Extreme Overfitting (High Variance)}: artificially high training performance ($R^2_{\text{train}}=0.850$) coexists with out-of-sample collapse (RSE ratio $>1.5$, $R^2$ loss $>0.20$). With $p=15$ predictors and only $n_{\text{train}}=96$ observations (6.4 observations per variable, far below a 15--20 safety threshold), the model memorized idiosyncratic noise rather than generalizable patterns.
\item Correction plan: (1) Ridge/Lasso regularization; (2) variable selection (recursive elimination or VIF $>5$) to retain 4--6 high-signal predictors; (3) nested cross-validation (inner hyperparameter loop and outer evaluation loop) to prevent information leakage; (4) log or Box--Cox transformations and robust regression to reduce the impact of high-leverage outliers.
\end{enumerate}
\end{solproblema}

% Analyze
\begin{solproblema}[prob:cb39746]
\begin{enumerate}
\item Model A: $\text{AIC}_A=100(1.5041)+2(3)=156.41$; $\text{BIC}_A=150.41+3(4.6052)=164.23$. Model B: $\text{AIC}_B=100(1.3350)+2(5)=143.50$; $\text{BIC}_B=133.50+5(4.6052)=156.53$. Model C: $\text{AIC}_C=100(1.2528)+2(8)=141.28$; $\text{BIC}_C=125.28+8(4.6052)=162.12$.
\item The smallest value is selected. AIC selects \textbf{Model C} (141.28); BIC selects \textbf{Model B} (156.53).
\item AIC penalizes each additional parameter by a fixed $+2$, whereas BIC uses $+\ln(n)$, exceeding 2 as soon as $n\ge8$. For $n=100$, each extra parameter costs $+4.61$ under BIC versus $+2.00$ under AIC. In Big Data, AIC tends to overselect variables with tiny effects, while BIC is asymptotically consistent and favors parsimonious models.
\end{enumerate}
\end{solproblema}

% Evaluate
\begin{solproblema}[prob:2b2e597]
\begin{enumerate}
\item Model B has greater raw explanatory power (0.800 versus 0.780). Coefficients of variation are $CV_A=0.040/0.780\times100\approx5.13\%$ and $CV_B=0.060/0.800\times100\approx7.50\%$; \textbf{Model A is more stable}.
\item With $n=150$, $R^2_{\text{adj},A}=1-[(149/146)(0.220)]\approx0.7755$ and $R^2_{\text{adj},B}=1-[(149/144)(0.200)]\approx0.7931$.
\item Model B's adjusted gain is only 1.76 percentage points, at the cost of 66.7\% more variables and substantially greater fold instability. Under Occam's razor, \textbf{recommend Model A}: its simplicity, stability, and nearly identical predictive precision protect against overfitting without sacrificing relevant power.
\end{enumerate}
\end{solproblema}

% Create
\begin{solproblema}[prob:9689172]
\begin{align}
\Delta R^2=R^2_{\text{train}}-R^2_{\text{test}}=0.82-0.79=0.03\text{ (3.0 percentage points)},
\end{align}
\begin{align}
\text{Relative RSE increase}=\frac{9.6-8.5}{8.5}\times100\approx12.94\%.
\end{align}
Since $\Delta R^2=0.03\le0.05$ and the relative RSE increase (12.94\%) is below 15\%, the model has \textbf{mild/admissible overfitting}, consistent with normal sampling fluctuation in a test sample of 150 customers. It generalizes reasonably to new customers and can be deployed, with periodic monitoring as subscriber behavior changes.
\end{solproblema}
